%% kapa: exam=1 round=28 year=2017 period=1 form=A subject=civil_law subject_ko=민법 seq=25 number=25
%% source: Q-Net 감정평가사 기출문제 (cst00309) artlSeq=5199603 file=extracted/1008/5199603/2017년 제28회 감정평가사 1차시험 문제지/1교시 A형.pdf
%% conversion: kordoc 4.24.51 -> tools/gen1.py (auto); review=proofread-sonnet; flags=none
\begin{question}{25}
선의취득에 관한 설명으로 옳지 않은 것은? (다툼이 있으면 판례에 따름)
\begin{choices}
  \item 경매에 의하여 소유권을 취득한 매수인에게도 선의취득이 인정될 수 있다.
  \item 동산의 선의취득은 양도인이 무권리자라고 하는 점을 제외하고는 아무런 흠이 없는 거래행위이어야 성립한다.
  \item 연립주택의 입주권은 선의취득의 대상이 될 수 없다.
  \item 저당권은 선의취득할 수 없다.
  \item 현실의 인도를 받지 않아도 점유개정의 방법만으로 선의취득이 인정된다.
\end{choices}
\end{question}
\answer{5}
